\chapter{\MakeUppercase{Introduction}}
\justifying
\section{Background}
Devices are typically secured via passwords, tokens, or certificates. While these types of keys are useful and well understood, they do lead to typical weaknesses associated with those methodologies, such as account sharing/access, stolen keys, replaying an old key, and using a key that is an explicitly known thing. Research in security has also explored the potential to use a physical characteristic of a device as a way to verify the identity of that device. This idea exists in hardware security, biometrics, and various forms of literature related to the measurement and evaluation of behavioral identity. The focus of this project is very similar to the previous examples; can the noise generated from sensors (e.g. camera/microphone) on a device be processed into a consistent and repeatable signal so that a high degree of confidence can be achieved in associating an identity to a measurement device?

The project approaches this issue as one of similarity (as opposed to identity). This distinction is essential to the success of the project. The raw or learned embedding would not be considered a password, token, or permanent secret; it is a measurement device. It must be used with newly acquired data (measurements) through the same measurement path and within a close range of the template to be of value. As a result, the project does not suggest that the noise generated from measurement devices can be considered a secure cryptographic identity; it only suggests that the potential for the noise to act as a viable means of verification exists.

Signal processing machine learning, and application security are the foundation of the body of work. Sensor noise is always statistical in nature and dependent on the environment, so the system does not depend on the identicalness of the same measurement being taken twice; it instead extracts engineered statistical and spectral features from the measurement and maps them into a learned embedding space, where proximity is measured using cosine distance. The entire pipeline is wrapped in a web service and has a database stored persistence, challenge expiration, rate limits, and audit logging. Therefore, the final outcome is a capstone project that is both model-driven and system-driven.

\subsection{Overview of the Problem Domain}
While every camera sensor and microphone path introduces small discrepancies in acquiring raw data, all cameras exhibit a range of qualities from fixed pattern noise to quantization effects as well as differences in noise caused by the device pipeline. Similarly, microphones will expose various sources of noise due to their electrical components and gain structure, as well as ambient variances that will affect how microphones work together in different settings. While none of these characteristics are guaranteed to be consistent or stable across various settings, they can therefore produce different measurable structures that may clearly differentiate samples from identical devices and samples from non-identical devices with a high degree of accuracy.

In the context of research, this presents an interesting problem. The signal is noisy, the classes are very similar, and the security environment must be sufficiently considered with respect to possible failure modes. Therefore, the challenge cannot be properly addressed by simply using a one-shot classifier. The resolution of this issue requires sampling over time, stable features, an accurate threshold calibration process, and clear assumptions regarding the intent of potential attackers.

\subsection{Current Systems and Existing Solutions}
Most deployed systems often use static credentials or hardware-bound cryptographic modules that provide effective performance for their intended uses, but for somewhat different issues (explicit identity proof vs. latent physical similarity). Though closer to the solution discussed in this thesis than other existing systems (like biometrics, device fingerprinting, and behavioral analytics), they generally need to operate with a larger scope of additional signals. QNA-Auth narrows its focus to both camera and microphone noise to determine if a relatively simple but total software system could utilize these signals to accomplish prototype-level verification.

\subsection{Limitations of Existing Approaches in the Project Context}
The following three considerations shaped the current project. First, the fact that methods based exclusively on secret signatures do not verify that a device's measurement path can help identify the author of it or not. Second, a large number of academic prototypes never progress beyond paper to become useful applications. Third, several experimental set-ups overstate their level of performance by collapsing probabilistic measurements of uncertainty into a ‘pass’ or ‘fail’ binary outcome. The current project addresses these shortcomings by constructing an end-to-end stack and providing method support for all documents generated by the military personnel who conducted the study as well as confidence bands and source-profile guards for evaluating the confidence of the measurements produced.

\section{Motivation}
This capstone was designed to meet both academic and practical objectives. From an academic standpoint, this capstone will synthesize signal processing, metric learning assessments, calibration of thresholds, and design of a secure system into one unified thesis. Practically, this capstone project will produce a live demonstration of a system that implements backend APIs, front-end displays/screens, an enrollment and authentication process/flow and review-ready artifacts. As a result, the capstone is beneficial for not only being an exercise in researching but demonstrating a high level of software engineering rigor.

There is also a methodological rationale behind this design. The primary purpose of a capstone project is to avoid claiming capabilities that the capstone will not be able to substantiate. The QNA-Auth project repository illustrates this approach through its repeated references to the system as being capable of “high-confidence device matching” but not “full identity proofing." The reasons for maintaining this usage of language are not merely rhetorical caution; they are an integral part of the technical correctness of the project.

\section{Scope of the Project}
This project aims to develop a multi-source enrollment solution, a runtime authentication engine, a challenge generation tool, a challenge verification engine, a device listing function, a device metadata retrieval function, and a device deletion function. The active sources for runtime evaluation are the camera and the microphone. Training and evaluation scripts, a review brief generator, diagnostic tools, and a React-based interface will also be part of the overall scope for the project.

There is no effort to provide a solution for large scale deployment of calibration, hardware attestation, universal cross-device uniqueness, and adaptive spoofing resistance. This project does not provide a substitute for cryptographic credentials. The appropriate interpretation of the intent of the deployment is an auxiliary verification prototype.

\section{Need for the Proposed System}
This system proposal meets the practical need between the "proof-of-concept" type signal analysis and usable, secure software, by converting sensor-noise experimentation into an organized service. Four distinct areas demonstrate the need for this system:
\begin{itemize}
	\item an enrollment workflow that is formalised so that it is not left to chance,
	\item a database supported method of storing device metadata and challenges,
	\item enhanced security to prevent replay using nonce binding; and
	\item reviewer-safe evaluation and demonstration narrative that are based on a real artifact from a repository.
\end{itemize}

\section{Organization of the Thesis}
This thesis consists of 8 sections. The purpose and objectives of the project, as well as how to fill in any gaps in knowledge, are covered in Chapter 2. Technical specifications and requirements are discussed in Chapter 3. The chapter describes how the system was designed. Chapter 4 describes the methodology and approach to testing., Chapter 5 Explain how the system is designed. It also describes how to build (implement) the system. The results and restrictions will be examined further in Chapter 6. Finally, Chapter 7 will give an overview of the work completed and provide recommendations for next steps.